\section{mObywatel (wersja 4.64.0 \,(4899))}
\label{sec:mobywatel}

\subsection*{Okno i wolumen}
Okno eksperymentu: \textbf{2025-08-24 13{:}23{:}03–13{:}26{:}00} ($\approx$ \SI{2.9}{\minute}). 
Zarejestrowano \textbf{13} przepływy, wolumen $\approx$ \textbf{0.21 MiB} 
(\textbf{RX $\approx$ 0.15 MiB}, \textbf{TX $\approx$ 0.06 MiB}).

\paragraph{Przebieg badania:} 
Scenariusz obejmował około \textbf{15 minut interakcji} z aplikacją, w tym: 
\emph{nasłuch tła (2 min), logowanie (biometria + hasło), wyświetlenie numeru PESEL, wgląd w dane wrażliwe (dowód, prawo jazdy, legitymacja studencka, dane pojazdu), nawigacja po głównych ekranach (offline i online), zamknięcie aplikacji oraz nasłuch tła (ok. 10 min)}.

\subsection{Wolumen i liczność wg protokołu}
\begin{table}[H]
  \centering
  \caption{Podsumowanie wg protokołu (mObywatel, ta sesja)}
  \label{tab:mobywatel-proto}
  \begin{tabular}{lrrrr}
    \toprule
    \textbf{Protokół} & \textbf{Flows} & \textbf{MiB (total)} & \textbf{B wysłane} & \textbf{B odebrane} \\
    \midrule
    HTTPS & 3  & 0.18 & 51{,}439 & 134{,}763 \\
    TLS   & 4  & 0.02 & 13{,}802 & 10{,}017 \\
    HTTP  & 1  & 0.01 & 630      & 8{,}157 \\
    DNS   & 5  & 0.00 & 327      & 407 \\
    \midrule
    \textbf{Razem} & \textbf{13} & \textbf{0.21} & \textbf{66{,}198} & \textbf{153{,}344} \\
    \bottomrule
  \end{tabular}
\end{table}

\paragraph{Obserwacja:} 
Dominuje \textbf{HTTPS/TLS}; pojedynczy ślad \textbf{HTTP/80} dotyczy hosta Google (prawdopodobnie kontrola łączności/zasób systemowy) i nie wskazuje na przesyłanie treści aplikacyjnych w badanym oknie czasowym.

\subsection{Największe cele (IP/port) wg wolumenu}
\begin{table}[H]
  \centering
  \caption{Największe kierunki ruchu (mObywatel, ta sesja)}
  \label{tab:mobywatel-topdst}
  \begin{tabular}{lllrr}
    \toprule
    \textbf{IP docelowe} & \textbf{Port} & \textbf{Proto} & \textbf{Flows} & \textbf{MiB} \\
    \midrule
    185.41.93.206  & 443  & HTTPS & 3 & 0.18 \\
    185.32.48.239  & 3000 & TLS   & 4 & 0.02 \\
    142.250.203.195 & 80  & HTTP  & 1 & 0.01 \\
    10.215.173.2   & 53   & DNS   & 5 & 0.00 \\
    \bottomrule
  \end{tabular}
\end{table}

\paragraph{Obserwacja:}
Ruch koncentruje się na 443/TCP do usług państwowych; widoczny jest także kanał TLS na nietypowym porcie \textbf{3000} (telemetria/błędy).

\subsection{Najczęstsze hosty (liczba przepływów)}
\begin{table}[H]
  \centering
  \caption{Top hosty (mObywatel, liczba przepływów)}
  \label{tab:mobywatel-hosts}
  \begin{tabular}{lr}
    \toprule
    \textbf{Host} & \textbf{Flows} \\
    \midrule
    \texttt{api.mobywatel.gov.pl} & 6 \\
    \texttt{sentry.puch.coi.gov.pl} & 5 \\
    \texttt{www.gstatic.com} & 2 \\
    \bottomrule
  \end{tabular}
\end{table}

\paragraph{Obserwacja:}
\texttt{api.mobywatel.gov.pl} dominuje; \texttt{sentry.puch.coi.gov.pl} wskazuje na wewnętrzny system raportowania błędów; \texttt{www.gstatic.com} występuje pomocniczo (zasoby/kontrola łączności).

\subsection{Obserwacje jakościowe}
\begin{itemize}
  \item \textbf{Szyfrowanie transportowe:} właściwa komunikacja aplikacji po TLS/443; brak treści w badanym oknie czasowym.
  \item \textbf{Brak QUIC/HTTP/3} w zarejestrowanym wycinku — kanał klasyczny TLS/TCP.
  \item \textbf{Telemetria:} obecność endpointu \texttt{sentry.puch.coi.gov.pl} (raportowanie błędów); wolumen niewielki.
  \item \textbf{HTTP/80 do Google:} pojedynczy flow (Google), typowy dla sprawdzeń łączności; brak wskazań na transfer danych aplikacji w badanym oknie czasowym.
\end{itemize}

\subsection{Wyniki z Exodus Privacy  \cite{ExodusPrivacy}}
Raport dla \texttt{pl.nask.mobywatel} (wersja 4.64.0 \,(4899), źródło Google Play) wskazuje \textbf{1 tracker} (\emph{Sentry} — crash reporting) oraz \textbf{20 uprawnień}. 

\textbf{Uprawnienia — przykłady kluczowe:}
\begin{itemize}
  \item \textbf{Lokalizacja:} \texttt{ACCESS\_COARSE\_LOCATION}, \texttt{ACCESS\_FINE\_LOCATION},
  \item \textbf{Multimedia/kamera/pamięć:} \texttt{CAMERA}, \texttt{READ\_EXTERNAL\_STORAGE}, \texttt{WRITE\_EXTERNAL\_STORAGE},
  \item \textbf{Biometria:} \texttt{USE\_BIOMETRIC}, \texttt{USE\_FINGERPRINT},
  \item \textbf{Telefonia/tożsamość:} \texttt{READ\_PHONE\_STATE},
  \item \textbf{Sieć i usługi:} \texttt{INTERNET}, \texttt{ACCESS\_NETWORK\_STATE}, \texttt{ACCESS\_WIFI\_STATE}, \texttt{FOREGROUND\_SERVICE}, \texttt{RECEIVE\_BOOT\_COMPLETED}, \texttt{WAKE\_LOCK}, \texttt{NFC}, \texttt{POST\_NOTIFICATIONS}, \texttt{GET\_TASKS},
  \item \textbf{Powiadomienia/FCM:} \texttt{RECEIVE}, \texttt{pl.nask.mobywatel.DYNAMIC\_RECEIVER\_NOT\_EXPORTED\_PERMISSION}.
\end{itemize}

\emph{Spostrzeżenia:} brak większej ilości trackerów potwierdza szczególną wagę, jaką aplikacja przykłada do prywatności. Zestaw uprawnień jest szeroki, ale adekwatny do funkcji (obsługa dokumentów tożsamości, NFC, biometria).

\subsection{Subiektywne spostrzeżenia}
\paragraph{Obserwowalność treści.} 
Ze względu na TLS treści zapytań/odpowiedzi nie są dostępne; nie widać przesyłania danych wrażliwych w badanym oknie czasowym.

\paragraph{Praktyki bezpieczeństwa.}
W analizowanym wycinku widać konsekwentne użycie HTTPS/TLS do usług \texttt{*.gov.pl}, brak nietypowych punktów końcowych poza kanałem telemetrycznym (TLS/3000) i brak szerokich odwołań do domen trzecich. To zachowanie spójne z wysokim standardem bezpieczeństwa na warstwie transportowej. Pełna ocena (np. pinning certyfikatów, mechanizmy anti-tamper/anti-root, magazyn kluczy, ochrona ekranu) wymaga analizy APK oraz testów dynamicznych.

\subsection{Podsumowanie}
W badanej sesji mObywatel wykazuje:
\begin{enumerate}
\item dominację TLS/HTTPS; wolumen ($\approx 0.21$ MiB), \textbf{13} przepływów w $\sim$\SI{2.9}{\minute},
\item koncentrację ruchu na \texttt{api.mobywatel.gov.pl} oraz wewnętrzną telemetrię \texttt{sentry.puch.coi.gov.pl},
\item pojedynczy flow HTTP/80 do zasobów Google (nie dotyczy danych wrażliwych aplikacji).
\end{enumerate}
